% Pro A. Cluentio (pro-cluentio) --- English translation
% Translated by Claude (Anthropic) from the Latin (Perseus, Clark text).
% Style guide: STYLE.md.
%
% The longest of Cicero's forensic speeches (202 sections),
% delivered in 66 BC. The defence of Aulus Cluentius Habitus on
% charges of poisoning his stepfather Statius Albius Oppianicus,
% complicated by Cluentius's own earlier prosecution of Oppianicus
% for poisoning -- a notoriously corrupt trial that has weighed
% on Cluentius's reputation through eight years of public hostility.

\ciceroSection{1} I have noticed, gentlemen, that all the prosecutor's speech is divided into two parts, of which the one seemed to me to lean and to trust greatly on the now grown-old hatred of the Junian trial; the other only for the sake of custom, timidly and distrustfully, to touch the reasoning of the charges of poisoning, on which matter this inquiry has been set up by the law. So I am resolved to keep this same division of hatred and charges in my defence in such a way that all may understand that I have neither wished to escape by being silent nor to obscure by speaking.

\ciceroSection{2} But when I consider in what manner I must labour in each, the one part --- which is proper to your trial and to the lawful inquiry on poisoning --- seems to me to be very brief and not to require great contention in speaking. The other, however, which is far removed from the trial, which is more fitted to seditiously roused public meetings than to calm and moderate trials --- I see how much difficulty in pleading and how much labour it will have.

\ciceroSection{3} But in this difficulty this thing, gentlemen, comforts me: that you are wont to hear charges so as to seek of the speaker the dissolving of all of them; that you reckon that you ought not to grant the defendant more for safety than the defender has been able by purging the charges and proving by speaking. About hatred, however, you ought so to debate among yourselves that you should consider not what is said by us but what ought to be said. For in the charges Aulus Cluentius's own peril is at issue; in hatred the common cause. Wherefore the one part of the case we shall plead so that we should teach you; the other so that we should beg. In the one your diligence is to be joined to us; in the other your faith is to be implored. For there is no one who can withstand hatred without your aid and that of such men.

\ciceroSection{4} Indeed, as far as it concerns me, I do not know whither to turn. Shall I deny that there was that infamy of a corrupted trial? Shall I deny that this matter was agitated in public meetings, tossed about in trials, recalled in the senate? Shall I tear from men's minds so great an opinion, so deeply set, so old? It is not of our wit; it is, gentlemen, of your help, in this calamitous fame, as if in some most ruinous flame and in a common fire, to come to the help of this innocence.

\ciceroSection{5} For just as in other places truth has too little of stay and too little of strength, so in this place false hatred ought to be weak. Let it lord it in public meetings; let it lie low in trials. Let it have weight in the opinions and talk of the unskilled; let it be rejected by the wits of the prudent. Let it have vehement sudden onsets; let it grow old with the interval of space and the case learned. Finally let that definition of fair trials which has been handed to us by our ancestors be kept: that in trials both fault should be punished without hatred and hatred set down without fault.

\ciceroSection{6} Wherefore I ask of you, gentlemen, before I begin to speak about the case itself, these things: first, that which is most fair --- that you should bring no prejudgement here. For we shall lose not only the authority but even the name of judges, if we shall not judge here from the cases themselves --- if we shall bring our judgements ready-made from home to the cases. Next, if you have already taken any opinion in your minds, if reasoning shall pluck it out, if speech shall shake it, if at last truth shall wrest it free, do not resist; let it go from your minds either willingly or with calm. But when I shall speak and dissolve about each thing, do not yourselves silently set against your thought what is contrary, but await the end and suffer me to keep my own order of speaking. When I shall have closed, then if anything shall have been passed over, ask of it in your minds.

\ciceroSection{7} I, gentlemen, easily understand that I am coming to that case which has now for eight continuous years been heard from the contrary side, and is by men's own opinion silently almost convicted and condemned. But if some god shall have won me your goodwill for hearing me, I shall surely bring it to pass that you understand that nothing is so to be feared by a man as hatred; that nothing is so to be desired by an innocent man, with hatred undertaken, as a fair trial, in which alone at last some end and outcome of false infamy is found. Wherefore great hope holds me, that, if I shall be able to set forth what is in the case and to attain all things by speaking, this place and your seating, which they thought would be dreadful and fearful for Aulus Cluentius, shall at last be the harbour and refuge of his wretched and much-tossed fortune.

\ciceroSection{8} Although there are very many things which seem to me ought to be said, before I speak about the case, of the common perils of hatred, yet, lest your expectation be longer held in suspense by my speech, I shall come to the charge with that prayer, gentlemen, which I see I must use the more often: that you should hear me as if at this time this case were spoken first, as it is spoken --- not as if it had often been spoken and never proved. For on this day for the first time the means of dissolving that old charge has been given. Before this time error in this case and hatred has been at large. Wherefore, while I answer briefly and clearly to the prosecution of many years, I ask, gentlemen, that you should hear me, as you have set yourselves to do, kindly and attentively.

\ciceroSection{9} Aulus Cluentius is said to have corrupted by money the trial, that he might condemn his innocent enemy, Statius Albius. I shall show, gentlemen, first, since the head of that atrocity and hatred was that an innocent man was caught by money, that no one was ever called into court on greater charges, by graver witnesses; next, that such prejudgements have been made about him by the very judges by whom he was condemned that not only by the same but not even by any others could he in any way be acquitted. When I have shown these things, then I shall show what I see is most asked: that that trial was tried by money not by Cluentius, but against Cluentius. And I shall make you understand in that whole case what the matter itself bore, what error feigned, what hatred forged.

\ciceroSection{10} First then this is, from which it can be understood that Cluentius ought greatly to have trusted the case: that, trusting in the surest charges and witnesses, he came down to accuse. Here, gentlemen, I must briefly set forth those charges by which Albius was condemned. I ask of you, Oppianicus, to think that I am unwilling speaking of your father's case, drawn on by faith and the duty of defence. For if to you in the present I shall not be able, yet many means of giving satisfaction shall be given me at the time remaining; if I do not now satisfy Cluentius, the means of satisfying afterward shall not be given me. At the same time, who is there who ought to doubt to speak against the condemned and the dead for the unhurt and the living? When against him against whom it is said the condemnation has now taken away every danger of disgrace, and death even of grief; while he for whom we are speaking can take in nothing of offence without the bitterest sense of mind and trouble and the highest disgrace and foulness of life.

\ciceroSection{11} And that you may understand that Cluentius brought the name of Oppianicus into court not by a prosecutor's mind, not led on by any display or glory, but by unspeakable wrongs, by daily plots, with the peril of life set before his eyes --- I shall seek a beginning a little further off for showing the matter; which I ask, gentlemen, that you should not bear heavily. For with the beginnings learned, you will much more easily understand the ends. Aulus Cluentius Habitus, the father of this man, gentlemen, was not only of the town of Larinum (from which he was), but also of that region and neighbourhood, in virtue, standing, nobility, the chief. He, when he had died in the consulship of Sulla and Pompeius, left this man fifteen years of age, and a daughter, grown and marriageable, who in a short time after her father's death married Aulus Aurius Melinus, her cousin, a young man among the chief (as he was then reckoned) honourable and noble among his own. When this marriage was full of dignity,

\ciceroSection{12} full of concord, suddenly the unspeakable lust of a fierce woman arose, joined not only with disgrace but with crime. For Sassia, the mother of this Habitus --- for "mother" by me in all the case (although she shows hostile hatred and cruelty against him), "mother," I say, she shall be called; nor shall she ever so hear from me of her own crime and monstrousness as to lose the name of nature; for the more loving and more indulgent the very name of mother is, the more by greater hatred shall you reckon worthy the singular crime of that mother who for many years now, and now most of all, has wished her son killed --- this Habitus's mother, then, was caught with love (against what was right) for that young Melinus, her son-in-law. And not even that for long. In whatever way she could, she held herself in that lust. Then she began so to blaze with frenzy, was so kindled, was so carried by lust, that neither shame nor piety, nor the family's stain, nor men's standing, nor her son's grief, nor her daughter's mourning could call her back from that desire.

\ciceroSection{13} The young man's mind, not yet firm in counsel and reasoning, she lured by all those things by which that age can be caught and softened. The daughter, who was not only afflicted by that common feminine grief at such a husband's wrongs, but could not bear her mother's unspeakable taking-of-her-husband (about which she reckoned she could not even complain without crime), wished the rest to be ignorant of so great an ill of hers. In the hands and lap of this most loving brother of hers she was wasting away with mourning and tears.

\ciceroSection{14} But behold, a sudden divorce, which seemed to be the comfort of all ills! Cluentia leaves Melinus --- not unwilling, given such great wrongs; not gladly, given that he was her husband. Then indeed that distinguished and famous mother began openly to leap with joy, to triumph with gladness --- the conqueror of her daughter, not of her lust. She was unwilling that her standing should longer be hurt by hidden suspicions. The bridal bed which two years before she had spread for her own daughter as bride, in the same house, with her daughter cast out and driven away, she ordered to be adorned and spread for herself. The mother-in-law marries her son-in-law, with no auspices, no witnesses, with the dreadful omens of all.

\ciceroSection{15} O incredible crime of a woman, and (save this one) unheard in all life! O unbridled and untamed lust! O singular audacity! Did she not fear, if not the power of the gods and the fame of men, then that very night and those wedding torches, not the threshold of the chamber, not her daughter's bed, finally not the very walls, witnesses of the earlier marriage? She broke through and laid prostrate all things by greed and frenzy. Lust overcame shame, audacity fear, frenzy reason.

\ciceroSection{16} The son bore this common disgrace of family, kinship, and name heavily. But his trouble was being increased by the daily complaints and ceaseless weeping of his sister. He resolved nevertheless that nothing graver in such great wrongs and so great a crime of his mother had to be done by him than that he should not deal with her, lest those things which he could not see without the highest grief of mind, by his dealing with his mother, he should be reckoned not only to see but even to approve by his own judgement.

\ciceroSection{17} What was the beginning to him of the strife with his mother, you have heard. That this had to do with the case, when you have learned the rest, you will understand. For it does not escape me that, of whatever sort a mother be, yet in the trial of her son it is scarcely fitting that the foulness of a parent should be spoken. I should not be a fit man for any case, gentlemen, if this --- which is set and fixed in the common feelings of men and in nature itself --- I, who am drawn to defending men's perils, did not see. I easily understand that men ought not only to be silent about the wrongs of parents but even to bear them with calm mind. But I think those things which can be borne must be borne; those which can be silent, must be silent.

\ciceroSection{18} Aulus Cluentius has seen no calamity in his life, has come to no peril of death, has feared no ill, that has not all been forged and put forth from his mother. Who at this time would be silent about all things, and would suffer those things, if she could not by forgetfulness, yet by her own silence to be covered. But indeed it is so done that she can in no way at all be silent. For this very trial, this peril, that prosecution, all that supply of witnesses which is to be, was at the start adorned by his mother; by his mother at this time it is being arrayed and prepared with all her resources and supplies. She herself lately flew from Larinum to Rome for the sake of crushing him. The audacious, wealthy, cruel woman is at hand. She sets up the prosecutors, she arrays the witnesses; she rejoices in his squalor and filth; she desires his ruin; she longs to pour out all her own blood, provided only that she sees his poured out before. Unless you shall see all these things plainly in the case, reckon her named by us rashly. But if they shall be both open and unspeakable, you ought to forgive Cluentius for suffering these things to be said by me; you ought not forgive me, if I were silent.

\ciceroSection{19} Now I shall set out summarily on what charges Oppianicus was condemned, that you may see through both Aulus Cluentius's constancy and the reasoning of the prosecution. And first I shall show what was the cause of accusing, that you may see Aulus Cluentius did this very thing compelled by force and necessity.

\ciceroSection{20} When he had caught manifestly the poison which Oppianicus, his mother's husband, had prepared for him, and the matter was held not by guess but by eyes and hands, and there could be no doubt at all in the case, he accused Oppianicus. With what constancy and with what diligence I shall say afterwards. Now I have wished you to know this: that there was for him no other cause of accusing save that he should shun the peril of life set before him and the daily snares laid for his head, by this one method. And that you may understand Oppianicus was prosecuted on such charges that neither the prosecutor ought to have feared nor the defendant to have hoped, I shall set forth a few of the charges of that trial. Which learned, none of you will wonder that he, distrusting his affairs, fled to Staienus and to money.

\ciceroSection{21} There was a certain Dinaea, of Larinum, mother-in-law of Oppianicus, who had three sons --- Marcus and Numerius Aurius and Gnaeus Magius --- and a daughter Magia married to Oppianicus. Marcus Aurius, a young man, was captured at Asculum in the Italic war, fell into the hands of Quintus Sergius the senator (he who was condemned among the assassins), and was at his house in the slave-prison. But Numerius Aurius, his brother, died and left Gnaeus Magius his brother as heir. Afterwards Magia, the wife of Oppianicus, died. Last, the only son left of Dinaea, Gnaeus Magius, died. He made heir that young Oppianicus, his sister's son, and ordered him to share with Dinaea his mother. Meanwhile an informer came to Dinaea --- not obscure or uncertain --- to tell her that her son Marcus Aurius was alive and was in the Gallic field in slavery.

\ciceroSection{22} The woman, with her children lost, when the hope of recovering one son was shown to her, called together all her kinsmen and her son's connections, and weeping asked of them that they should take up the business, track down the young man, restore to her the son whom yet alone of many fortune had wished to leave her. When she had begun to do these things, she was overcome by disease. So she made a will of such kind, that to that son she bequeathed 400,000 sesterces, and instituted as heir that same Oppianicus, her grandson. And in those few days she died. Yet her kinsmen, as they had set themselves while Dinaea lived, so when she had died, set out to track down Marcus Aurius into the Gallic field with that same informer.

\ciceroSection{23} Meanwhile Oppianicus, as he was --- as you will find from many things --- of singular crime and audacity, through a certain Gallicanus, his familiar, first corrupted that informer with money. Then he saw to the taking off and killing of Marcus Aurius himself, with no great loss made. But those who had set out to track down and recover the kinsman send letters to Larinum to the Aurii, the kinsmen of that young man and his connections: that the method of tracking him down was hard for them, because they understood that the informer had been corrupted by Oppianicus. Which letters Aulus Aurius, a brave and resourceful man, noble at home and the very near kinsman of that Marcus Aurius, openly in the forum, with many hearing, with Oppianicus present, reads out; and with the clearest voice testifies that he will lodge the name of Oppianicus, if he should learn that Marcus Aurius had been killed.

\ciceroSection{24} Meanwhile in a short time those who had set out to the Gallic field return to Larinum. They announce that Marcus Aurius has been killed. The minds not only of the kinsmen but even of all the Larinates are moved, by hatred of Oppianicus and by pity of that young man. And so when Aulus Aurius, who before had given notice, had begun to pursue the man with cries and threats, he fled from Larinum and betook himself to the camp of the most distinguished man, Quintus Metellus.

\ciceroSection{25} But after that flight, the witness of crime and conscience, he never dared to commit himself to trials, never to laws, never unarmed to his enemies. But through that force and victory of Lucius Sulla, with armed men he flew up to Larinum in the highest fear of all. He took away the four-men whom the townsmen had made; he said that he had been made by Sulla, and three others besides, and that by the same it had been ordered him that Aulus Aurius (who had shown to him the lodging of his name and the peril of his life), and the second Aulus Aurius and his son Lucius, and Sextus Vibius (whom he was said to have used as go-between in corrupting that informer) should be proscribed and killed. So with these most cruelly killed, the rest were terrified by no middling fear of proscription and death from him. With these things made plain in the case and trial, who is there who could think that he could have been acquitted? And these are small things; learn the rest, that you may wonder that Oppianicus was not at any time condemned, but was unhurt for some time.

\ciceroSection{26} First see the audacity of the man. He longed to take Sassia in marriage, the mother of Habitus --- her, whose husband, Aulus Aurius, he had killed. Whether is he the more shameless who asked it, or she the more cruel, if she should marry, is hard to say. But yet learn the humanity and constancy of each.

\ciceroSection{27} Oppianicus asks that Sassia marry him, and contends greatly. But she does not wonder at his audacity, does not despise his shamelessness, finally does not dread that house of Oppianicus dripping with the blood of her husband. But because he had three sons, she answered that on this account she shrank from that marriage. Oppianicus, who had longed for Sassia's money, thought a remedy must be sought at home for that delay which was being brought to the marriage. For when he had by Novia an infant son, but his other son (born by Papia at Teanum, which is eighteen miles from Larinum) was being brought up at his mother's, he sent for the boy from Teanum suddenly without cause, which he had not been wont to do save on public games or festal days. The wretched mother, suspecting nothing of ill, sends him. He, when he had pretended that he was setting out for Tarentum, on that very day the boy --- when at the eleventh hour he had been seen in public well --- before night was dead, and the next day before it was light was burned.

\ciceroSection{28} And to the mother this great mourning was reported by the rumour of men sooner than by any one out of Oppianicus's household. She, when she had heard at one time that not only her son but also the office of the funeral was snatched from her, came at once to Larinum unconscious, and there held the funeral over from the start, with her son already buried. Ten days had not yet passed when that other infant son is killed. So Sassia at once marries Oppianicus, with mind glad now and hope most well-confirmed; nor wonder, who saw herself softened not by wedding gifts but by the funerals of his sons. So, what others on account of children are wont to be more eager for money, he reckoned it pleasant to lose his children on account of money.

\ciceroSection{29} I see, gentlemen, that, by your humanity, with these so great crimes briefly shown by me, you are vehemently moved. With what spirit then do you suppose those men were, who about these matters not only had to hear but also to judge? You hear about him over whom you are not judges, about him whom you do not see, about him whom you can no longer hate, about him who has satisfied both nature and the laws --- whom the laws have punished with exile, nature with death. You hear not from an enemy, you hear without witnesses, you hear when those things which can be most copiously said are spoken briefly and in passing. Those men were hearing about him over whom they were sworn to give their votes, about him whose unspeakable and crime-laden face, present, they were looking at, about him whom all hated on account of his audacity, about him whom they reckoned worthy of every punishment. They were hearing from prosecutors; they were hearing the words of many witnesses; they were hearing about each thing while it was being spoken gravely and at length by Publius Cannutius, a most eloquent man.

\ciceroSection{30} And is there anyone who, when he has learned these things, can suspect that Oppianicus was crushed in trial and circumvented as innocent? In a heap, gentlemen, I shall now set out the rest, that I may come to those things which are nearer to this case and more closely joined. Hold, please, in memory that this is not set forth by me --- that I should accuse the dead Oppianicus --- but, since I wish to persuade you that the trial was not corrupted by him, that I should use this as the start and foundation of the defence: that Oppianicus, a most criminal and most guilty man, was condemned. Who when he himself had given the cup to his wife Cluentia, the aunt of this Habitus, suddenly she in the middle of the drinking cried out that she was dying with the greatest pain, nor did she live longer than she spoke. For in that very speech and shout she died. And to this sudden death and the voice of the dying, all besides which are wont to be the signs and traces of poison were on the body of the dead woman. And by the same poison he killed his brother Gaius Oppianicus.

\ciceroSection{31} Nor is this enough. Although in the very fratricide no crime seems to be passed over, yet, that he might come to this unspeakable deed, he prepared his approach by other crimes before. For when his brother's wife Auria was pregnant, and the birth was now thought to be near, he killed the woman with poison, that at the same time what had been conceived from his brother might be killed. After he attacked his brother, who, late, with that cup of death now drained, when he was crying out about his own and his wife's destruction and longed to change his will, in the very signifying of this will died. So he killed the woman, lest by her birth he should be shut out from his brother's inheritance. But his brother's children he deprived of life before they could receive this light from nature --- that all might understand that nothing could be closed, nothing sacred, against him from whose audacity the children of his brother not even the guarding of the mother's body could keep safe.

\ciceroSection{32} I hold in memory a Milesian woman, when I was in Asia, who, because she had taken money from the second heirs and procured an abortion for herself by drugs, was condemned of a capital matter; nor unjustly, who had taken away the hope of a parent, the memory of a name, the help of a line, the heir of a household, a destined citizen of the commonwealth. By how much is Oppianicus, in the same wrong, worthy of greater punishment! If indeed she, when she had brought force on her own body, tortured herself; but he brought about the same through another's body's death and torture. Others do not seem able to undertake many parricides in single men; Oppianicus has been found to kill more in one body.

\ciceroSection{33} So when his uncle had learned this custom and audacity of his --- the uncle of that young man Oppianicus, Gnaeus Magius --- and he, when he was afflicted with a grave disease and was making that man, his sister's son, his heir, with friends called in and his mother Dinaea present, asked his wife if she was pregnant. When she answered that she was, he asked of her that, when he was dead, at Dinaea's house --- who was then her mother-in-law --- until she should give birth, she should live, and bring diligence that what she had conceived she might be able to keep and to bring forth safe. So he bequeathed her by will a great sum of money from his son if any should be born; from the second heir nothing he bequeathed. What he suspected of Oppianicus you see;

\ciceroSection{34} what he judged is not obscure. For whose son he made his heir, him he did not name as guardian for his children. What Oppianicus did, learn, that you may understand that Magius did not look far ahead in mind as he was dying. The money which had been bequeathed to the woman from her son if any should be born, that, ready, Oppianicus paid the woman not owed --- if this is to be called payment of a legacy and not the wage of an abortion. By which price taken and by many gifts besides (which then were being read out from Oppianicus's tablets), the hope which she had in her belly entrusted by her husband, conquered by greed, she sold to the crime of Oppianicus. Nothing now seems able to be added to this wickedness;

\ciceroSection{35} mark the outcome. The woman who, by her husband's adjuration, in those ten months ought not to know any house save her mother-in-law's, in the fifth month after her husband's death married Oppianicus himself. Which marriage was not long-lasting; for they had been joined not by the dignity of marriage but by the partnership of crime.

\ciceroSection{36} What? That slaughter of Asuvius of Larinum, a wealthy young man --- how famous it was then, the matter recent; how celebrated by the talk of all! There was a certain Avillius of Larinum, of ruined wickedness and the highest want, endowed with a kind of art fitted to rousing the lusts of young men. Who, when he had plunged himself by flatteries and assents deeply into Asuvius's company, Oppianicus at once began to hope that, by Avillius as if by some engine moved up, he could capture Asuvius's youth and storm his ancestral fortunes. The plan was thought up at Larinum; the matter was carried to Rome. For they reckoned it could be more easily entered upon in solitude, and more conveniently finished in a crowd. Asuvius set out with Avillius for Rome. Oppianicus followed them on their tracks. How they lived at Rome, in what banquets, in what disgraces, with how great and how lavish expenses --- with not only Oppianicus aware but a fellow-banqueter and helper --- it is long for me to say, especially as I hasten to other things. Learn the outcome of this feigned familiarity.

\ciceroSection{37} When the young man was at the house of a certain little woman, and where he had spent the night was tarrying for the next day --- Avillius, as had been arranged, pretends to be sick and to wish to make his will. Oppianicus brings to him sealers who knew neither Asuvius nor Avillius, and calls that man Asuvius himself; with the will sealed in Asuvius's name they depart. Avillius at once recovers; but Asuvius in that short time, as if he were going to a little garden, was led out to certain sand-pits beyond the Esquiline gate and killed.

\ciceroSection{38} When he was already missed for one and a second day and was not being found in those places where by custom he was being looked for, and Oppianicus in the forum of the Larinates was repeatedly saying that lately he and his friends had sealed his will, the freedmen of Asuvius and not a few friends, because it was agreed that Avillius had been with him on the day on which Asuvius had last been seen, and had been seen by many, attack the man and set him before the feet of Quintus Manlius, who was then triumvir. And there at once, with no witness, no informer, terrified by the consciousness of the recent evil deed, he sets out everything (as it was said by me a little before) and confesses that Asuvius had been killed by him at the counsel of Oppianicus. Oppianicus is dragged out from the house where he was hiding by Manlius;

\ciceroSection{39} the informer Avillius is held openly on the other side. What now do you ask of the rest? Most of you knew Manlius. He had never thought of honour from boyhood, of the pursuits of virtue, of any fruit of good standing; but, from a petulant and dishonest buffoon, in the discords of the state had now come by the votes of the people to that column at which by the abuse of many he had often been carried. So then he settles with Oppianicus, takes money from him, leaves the case --- both undertaken and so plain. And then in Oppianicus's case this Asuvian charge was confirmed both by many witnesses and indeed by his very will, in which it was agreed that Oppianicus's name was first --- the man whom you say was wretched and innocent, circumvented by a false trial.

\ciceroSection{40} What? Your grandmother, Oppianicus, Dinaea, of whom you are heir --- did your father not manifestly kill? To whom, when he had brought that famous doctor of his, now well known and often victorious, through whom he had killed very many, the woman cries out that she does not at all wish to be cured by him by whose curing she had lost all her own. Then he suddenly attacks a certain Anconitan man, Lucius Clodius, a circuit pharmacist, who by chance had then come to Larinum, and settles with him at 2,000 sesterces (which is shown from his very tablets). Lucius Clodius, since he was hurrying, who had many circuits remaining, as soon as he was brought in, finished the matter; with the first dose he took the woman off and afterwards stayed not a moment of time at Larinum. The same Dinaea making her will,

\ciceroSection{41} when Oppianicus, who had been her son-in-law, had taken the tablets, with his finger he wiped out the legacies; and, when he had done this in many places, after her death, lest he be convicted by the erasures, he had the will copied onto other tablets and sealed it with adulterated seals. Many things I pass over on purpose; for I fear that these very things may seem too many. Yet you ought to reckon that he was like himself in the rest of the parts of his life. The decurions all judged that he had corrupted the public records of Larinum, the censorial ones. With him no one now had business, no one transacted any matter; no one out of so many kinsmen and connections ever wrote him as guardian for his children; no one judged him worthy of approach, no one of meeting, no one of speech, no one of banquet. All shrank away, all turned aside, all fled him as if some monstrous and ruinous beast and pestilence.

\ciceroSection{42} Yet this man so audacious, so unspeakable, so guilty Habitus would never have prosecuted, gentlemen, if he could have passed it over with his life safe. Oppianicus was an enemy to him, was so. But yet he was his stepfather. Cruel and hostile to him was his mother, but his mother. Finally nothing was so far from accusation as Cluentius, both by nature and by will and by the ordered reasoning of his life. But when this condition was set before him --- that either he should justly and dutifully accuse, or bitterly and unworthily die --- he preferred to accuse in whatever way he could, than to die in that manner.

\ciceroSection{43} And that you may see this is so, I shall set out for you Oppianicus's deed manifestly learned and caught --- from which at the same time you will understand that both for him to accuse and for him to be condemned was necessary. There were certain Martiales at Larinum, public ministers of Mars, consecrated to that god by the old institutions and religious sanctities of the Larinates. Of whom, since there was a great enough number, and likewise (as in Sicily there are very many Venusians) so they at Larinum were reckoned in Mars's household, suddenly Oppianicus began to defend that they were all free and Roman citizens. Heavily the decurions of Larinum and all the townsmen took it. So they asked Habitus that he should take up that case and defend it publicly. Habitus, although he had removed himself from every business of that kind, yet for the sake of his place, for the antiquity of his line, for that --- because he reckoned himself born not for his own advantages but also for those of his townsmen and other connections --- did not wish to fail the will of all the Larinates so great.

\ciceroSection{44} With the case undertaken and brought to Rome, daily great contentions between Habitus and Oppianicus were stirred up out of the zeal of each for the defence. Oppianicus himself was of monstrous and bitter nature; his frenzy was kindled by the mother of Habitus, hostile and inimical to her son. The Magni, however, thought it of their interest that he should be moved away from the Martiales' case. There lay underneath also another greater cause which moved most the mind of Oppianicus, a most greedy and most audacious man.

\ciceroSection{45} For Habitus up to the time of that trial had never made any will. For he could not bring it into his mind to leave anything of that sort to his mother, nor by will to pass over altogether the name of his parent. When Oppianicus knew this --- for it was not obscure --- he understood that, with Habitus dead, all his goods would come to his mother; who, with money increased by him, with greater gain, deprived of her son, with less peril could be killed. So, set on fire by these things, by what reasoning he tried to take away Habitus by poison, learn.

\ciceroSection{46} Gaius and Lucius Fabricius were twin brothers, of the town of Aletrium, men like one another both in form and in morals, but most unlike their townsmen --- of whom what splendour, how almost equal, how almost in all constant and moderate the reasoning of life is, none of you (in my opinion) is ignorant. With these Fabricii Oppianicus had always most familiarly used. Now you almost all know how great force the likeness of pursuits and nature has for joining friendships. When they so lived as to think no gain shameful, when from them every fraud, all snares and circumventions of young men were born, and they were known to all by vices and dishonesty --- zealously, as I said, Oppianicus had attached himself to their familiarity many years before. So then he resolved through Gaius

\ciceroSection{47} Fabricius (for Lucius had died) to prepare snares for Habitus. At that time Habitus was of weak health. He was using a not ignoble but a respected man as doctor, Cleophantus, whose slave Diogenes Fabricius began to solicit by hope and price for giving poison to Habitus. The slave, not uncunning and (as the matter itself made plain) thrifty and upright, did not despise Fabricius's speech; he took the matter to his master. Cleophantus then conferred with Habitus. Habitus at once shared it with Marcus Baebius, a senator, his most familiar friend; of whom, of what faith, of what prudence, of what diligence he was, I think you remember. It pleased him that Habitus should buy Diogenes from Cleophantus, that the matter could the more easily either be caught by his information or shown to be false. To be brief: Diogenes is bought; in a few days the poison is procured. Many good men, when they had come up secretly, the money sealed which was being given for that thing was caught in the hands of Scamander, freedman of the Fabricii.

